\chapter{Defect correction: the rate, never the answer}

Every other slot for a learned expert asks it to be part of the answer, and so
owes it a certificate.  Defect correction puts it somewhere else.  $\Phi$ is the
classical step, whose settled states are its fixed points.  $\Psi$ is any other
map on the same states: a learned operator, a coarse solver, anything cheap.
With $G_\Phi=I-\Phi$ and $G_\Psi=I-\Psi$, the next iterate $w'$ is defined by
\[
  G_\Psi(w')=G_\Psi(w)-G_\Phi(w),\qquad\text{that is,}\qquad
  w'-\Psi(w')=\Phi(w)-\Psi(w)
\]
(Stetter, \emph{Numer. Math.} 29, 1978).  The cheap map enters only through a
difference, and the classical map supplies the defect.

\section{T5: Theorem 1 and Corollary 1}

\paragraph{In plain words.}
If the iterates settle anywhere, they settle on a fixed point of the classical
map.  Nothing about the accuracy of the cheap map enters, only that both maps
are continuous at the limit.  A wrong cheap map can cost classical calls; it
cannot change what is returned.  And if the cheap map is constant, the
iteration is the classical march itself.

\begin{definition}[A defect-correction step]\label{def:dec-step}
  \lean{Atlas.DefectCorrection.IsStep, Atlas.DefectCorrection.innerMarch}
  \leanok
  Let $E$ be a normed space and $\Phi,\Psi:E\to E$.  The point $w'$ is a
  successor of $w$ if $w'-\Psi(w')=(w-\Psi(w))-(w-\Phi(w))$.  The \emph{inner
  march} that solves for it is $x^{(0)}=\Phi(w)$,
  $x^{(m+1)}=\Phi(w)+[\Psi(x^{(m)})-\Psi(w)]$.
\end{definition}

\begin{theorem}[T5, Theorem 1: consistency, whatever $\Psi$ is]\label{thm:T5}
  \lean{Atlas.DefectCorrection.limit_isFixedPt}
  \leanok
  \uses{def:dec-step}
  Let $w_{k+1}$ be a successor of $w_k$ for every $k$, and $w_k\to\bar w$.  If
  $\Phi$ and $\Psi$ are continuous at $\bar w$, then $\Phi(\bar w)=\bar w$.
\end{theorem}

\begin{proof}
  \uses{def:dec-step}
  $w_k-\Phi(w_k)=\bigl(w_k-\Psi(w_k)\bigr)-\bigl(w_{k+1}-\Psi(w_{k+1})\bigr)$.
  The right side tends to $(\bar w-\Psi(\bar w))-(\bar w-\Psi(\bar w))=0$ and the
  left to $\bar w-\Phi(\bar w)$.
\end{proof}

\begin{theorem}[T5, Theorem 1 with inexact inner solves]\label{thm:T5-inexact}
  \lean{Atlas.DefectCorrection.limit_isFixedPt_of_inexact}
  \leanok
  \uses{def:dec-step}
  Suppose instead that each step misses the defining equation by at most
  $\varepsilon_k$,
  \[
    \norm{\bigl(w_{k+1}-\Psi(w_{k+1})\bigr)
      -\bigl((w_k-\Psi(w_k))-(w_k-\Phi(w_k))\bigr)}\le\varepsilon_k ,
  \]
  with $\varepsilon_k\to0$.  If $w_k\to\bar w$ and $\Phi$, $\Psi$ are continuous
  at $\bar w$, then $\Phi(\bar w)=\bar w$.
\end{theorem}

\begin{proof}
  \uses{def:dec-step}
  As for Theorem~\ref{thm:T5}: the left side of the same identity now differs
  from the right by at most $\varepsilon_k\to0$.
\end{proof}

\begin{corollary}[T5, Corollary 1: the null element]\label{cor:T5-const}
  \lean{Atlas.DefectCorrection.step_of_const}
  \leanok
  \uses{def:dec-step}
  If $\Psi$ is constant and $w'$ is a successor of $w$, then $w'=\Phi(w)$.
\end{corollary}

\begin{proof}
  \uses{def:dec-step}
  $w'-c=(w-c)-(w-\Phi(w))$.
\end{proof}

\begin{lemma}[The inner march solves for the successor]\label{lem:T5-inner}
  \lean{Atlas.DefectCorrection.isStep_iff}
  \leanok
  \uses{def:dec-step}
  $x$ is a successor of $w$ if and only if $x=\Phi(w)+[\Psi(x)-\Psi(w)]$.
\end{lemma}

\begin{proof}
  \uses{def:dec-step}
  Rearrangement.
\end{proof}

\begin{corollary}[T5, Corollary 1 for the inner march]\label{cor:T5-inner-const}
  \lean{Atlas.DefectCorrection.innerMarch_of_const}
  \leanok
  \uses{def:dec-step}
  If $\Psi$ is constant, $x^{(m)}=\Phi(w)$ for every $m$.
\end{corollary}

\begin{proof}
  \uses{def:dec-step}
  Induction on $m$; the bracket is zero.
\end{proof}

\section{T5: Corollary 2, shrinking toward the null element}

\paragraph{In plain words.}
The iteration's rate is set by $J_\Psi^{-1}J_\Phi$, with $J_\Psi=I-D\Psi$ the
derivative of $G_\Psi$ at the settled state.  A cheap map that is nearly
singular on some mode makes the iteration diverge.  Replace $\Psi$ by
$(1-\alpha)\Psi$.  Its $J$ is $\alpha I+(1-\alpha)J_\Psi$, whose eigenvalues are
those of $J_\Psi$ moved toward $1$; wherever $J_\Psi$'s eigenvalue has
non-negative real part, the new one has real part at least $\alpha$.

\begin{lemma}[The shrunk Jacobian]\label{lem:T5-shrink-deriv}
  \lean{Atlas.DefectCorrection.shrink_hasFDerivAt}
  \leanok
  If $\Psi$ has Fr\'echet derivative $D$ at $w$, then
  $x\mapsto x-(1-\alpha)\Psi(x)$ has derivative
  $\alpha I+(1-\alpha)(I-D)$ at $w$.
\end{lemma}

\begin{proof}
  Linearity of the derivative, and
  $I-(1-\alpha)D=\alpha I+(1-\alpha)(I-D)$.
\end{proof}

\begin{theorem}[T5, Corollary 2: the eigenvalues]\label{thm:T5-shrink}
  \lean{Atlas.DefectCorrection.shrink_hasEigenvalue_iff}
  \leanok
  \uses{lem:T5-shrink-eig}
  Let $J$ be a linear map on a vector space over a field, and $\alpha\ne1$.
  Then $\nu$ is an eigenvalue of $\alpha I+(1-\alpha)J$ if and only if
  $\nu=\alpha+(1-\alpha)\mu$ for an eigenvalue $\mu$ of $J$.
\end{theorem}

\begin{proof}
  \uses{lem:T5-shrink-eig}
  Lemma~\ref{lem:T5-shrink-eig} gives one direction.  For the other,
  $(\alpha I+(1-\alpha)J)v=\nu v$ gives $Jv=\frac{\nu-\alpha}{1-\alpha}v$.
\end{proof}

\begin{lemma}[Eigenvalues move with the shrink]\label{lem:T5-shrink-eig}
  \lean{Atlas.DefectCorrection.shrink_hasEigenvalue}
  \leanok
  If $Jv=\mu v$ then $(\alpha I+(1-\alpha)J)v=(\alpha+(1-\alpha)\mu)v$.
\end{lemma}

\begin{proof}
  Direct.
\end{proof}

\begin{lemma}[The real part]\label{lem:T5-shrink-re}
  \lean{Atlas.DefectCorrection.shrink_re}
  \leanok
  For $0\le\alpha\le1$ and complex $\mu$ with $\operatorname{Re}\mu\ge0$,
  $\operatorname{Re}\bigl(\alpha+(1-\alpha)\mu\bigr)\ge\alpha$.
\end{lemma}

\begin{proof}
  $\operatorname{Re}(\alpha+(1-\alpha)\mu)=\alpha+(1-\alpha)\operatorname{Re}\mu$.
\end{proof}
